\subsection{Merge sort}

Merge sort resuelve el problema de ordenamiento (definido en la sección \ref{subsec:problema}).

La idea es dividir el arreglo a ordenar en dos mitades, ordenar cada mitad de forma recursiva (es decir, dividiendo cada una otra vez a la mitad, y así sucesivamente) y luego combinar las soluciones.

Haciendo explícitas las etapas:
\begin{itemize}
    \item \textbf{Dividir}: Dividir el arreglo en dos mitades.
    \item \textbf{Conquistar}: Ordenar cada mitad usando merge sort (o sea, aplicando estas etapas a cada mitad)
    \item \textbf{Combinar}: Combinar las soluciones de cada mitad, que son subarreglos ordenados.
\end{itemize}

Los primeros dos pasos son muy sencillos porque en algún punto al dividir por la mitad se obtienen arreglos de un elemento, que ya no hay que ordenar. Entonces, si encapsulamos el paso de combinar, este es el seudocódigo:
\begin{pseudocode}
función merge_sort(números: arreglo, pos_min: entero, pos_max: entero)
    if pos_min >= pos_max
        return números

    pos_mitad = (pos_min + pos_max) // 2
    merge_sort(números, pos_min, pos_mitad)
    merge_sort(números, pos_mitad+1, pos_max)

    combinar(números, pos_min, pos_mitad, pos_max)
\end{pseudocode}
Necesitamos tener los índices como parámetros porque indican cómo se va dividiendo el arreglo. La invocación principal, que resuelve el problema, requiere los índices extremos del arreglo: \inline{merge_sort(números, 0, números.longitud-1)}. Nótese que en esta función el arreglo realmente no se divide: las llamadas recursivas trabajan sobre el mismo arreglo con distintos índices, la división es conceptual.
% TODO: Explicar, o poner como referencia, a CLRS
% Explicar que acá el arreglo pasa por referencia, no por valor

Sorprendentemente, la mejor forma de combinar las soluciones es la forma más intuitiva. Si cada arreglo fuera una pila de cartas, el método natural sería comparar la primera de cada pila y tomar la menor; luego comparar la siguiente de ese arreglo con la primera del otro arreglo y tomar la menor (que es la que va después de la primera que tomé), y así sucesivamente. El seudocódigo es el siguiente:
\begin{pseudocode}
función combinar(números: arreglo, pos_min: entero, pos_mitad: entero, pos_max: entero)
    izquierda = |\textrm{copia de números desde}| pos_min |\textrm{hasta}| pos_mitad
    derecha = |\textrm{copia de números desde}| pos_mitad+1 |\textrm{hasta}| pos_max
    i = 0; j = 0

    |\br{ai}|while i < izquierda.longitud and j < derecha.longitud
        if izquierda[i] |\bxl{es}|<=|\bxr{es}| derecha[j]
            números[pos_min+i+j] = izquierda[i]
            i++
        else
            números[pos_min+i+j] = derecha[j]
            j++|\br{af}|

    |\br{ii}|while i < izquierda.longitud
        números[pos_min+i+j] = izquierda[i]
        i++|\br{if}|

    |\br{di}|while j < derecha.longitud
        números[pos_min+i+j] = derecha[j]
        j++|\br{df}|
\end{pseudocode}
\begin{annotations}
    \codebrace[xshift=7.8cm, width=4.5cm]{ai}{af}{Ambos subarreglos tienen elementos; se toma el menor y se añade a \texttt{números}.}
    \codebrace[xshift=6cm, width=7cm]{ii}{if}{Si quedan elementos en el arreglo de la izquierda, ya no quedan en el de la derecha. Se añade el excedente de la izquierda a \texttt{números}.}
    \codebrace[xshift=6cm, width=7cm]{di}{df}{Se añade el excedente de la derecha a \texttt{números}. Si este while corre, es porque el anterior no corrió y viceversa.}
    \codebox[notepos={(cbox@es.east)+(2cm,0.8cm)}, curve={bend left=25}, width=5cm]{es}{\texttt{<=} en lugar de \texttt{<} para estabilidad.}
\end{annotations}

Un algoritmo es \emph{estable} si mantiene el orden relativo de los elementos iguales. O sea, en caso de empate, deja primero al que estaba primero en el arreglo original. En la línea del comentario sobre estabilidad, se elige el elemento de la izquierda en caso de empate para que continúe ubicado antes que el de la derecha en la solución final.

% TODO: Análisis de complejidad? Así sea breve
% Ojo: para poder decir que es logarítmico formalmente, necesito el teorema maestro